\documentclass[11pt]{report}
\usepackage{amsmath, amssymb, amsthm, bbm}
\usepackage{graphicx}
\usepackage{hyperref}
\usepackage{tikz, subcaption}
\usepackage{enumitem}
\usepackage{geometry}
\usepackage{algorithm}
\usepackage{algpseudocode}
\usepackage[backend=biber,
    maxnames=10,
    minnames=1,
style=alphabetic]{biblatex}
\addbibresource{references.bib}
\geometry{margin=1in}

\newtheorem{proposition}{Proposition}
\newtheorem{theorem}{Theorem}
\newtheorem{lemma}{Lemma}
\newtheorem{corollary}{Corollary}
\newtheorem{definition}{Definition}[section] 

\title{Linear Programming of Graph Squares}


\author{
    Eddie Xu \\[2em]
    Submitted for the degree of Master of Science:\\
    Mathematics and Statistics
}

\date{
    \vspace{2em}
    Supervisor: Christopher Duffy \\[1em]
    \today
} 


\begin{document}

\maketitle


\tableofcontents

\newpage 

\chapter{Introduction} \label{chap::intro}

\section{Historical Background}

In 1949, Leon Festinger proposed a novel method in the field of social psychology for analyzing relationships among a number of people 
present \cite{festinger_bahoken_beauguitte_2023}. Classically this was represented using a sociogram, a visual tool that maps out the 
relationships between people, however Forsyth and Katz introduced the use of a matrix to represent the sociometric pattern, with rows and 
columns being individuals and containing $1$s for choices and blank diagonals. Festinger noticed that the action of squaring a 
matrix 
\[
A_{rc}^2 = A_{1c}A_{r1} + \dots + A_{nc}A_{rn} \footnote{$A_{rc}^2$ is more naturally written as $\displaystyle \sum_{k = 1}^n A_{ck}A_{kr}$
using the commutativity of multiplication $A_{1c}A_{r1} = A_{c1}A_{1r}$, which intuitively gives the notion of paths from c to r via intermediary k}
\]
represented the number of ``two-step'' connections, which can then be used to denote ``three-step'' connections between any two people  
\[
A_{rc}^3 = A_{1c}A_{r1}^2 + \dots + A_{nc}A_{rn}^2, \footnote{$A_{rc}^3$ is more naturally written as $\displaystyle \sum_{k=1}^n A_{ck}A^2_{kr}$,
which intuitively counts the number of paths from $c$ to $r$ that go through an intermediate vertex $k$, followed by any two-step path from 
$k$ to $r$.}
\]
such that higher steps can be easily generated. In this manner, higher powers of the adjacency matrix encode the number of walks of increasing 
length between pairs of individuals, allowing densely connected groups of people, known as cliques, to be identified. In particular, the 
cliques between three individuals can be done by looking at the diagonals of the cubed matrix $A^3$.  

This matrix $A$ can be considered as a form of adjacency matrix, and in the process of improving clique detection methods, 
\citeauthor{ross_1960_square} showed that complete graphs may be realized as squares of star trees by considering powers of such matrices 
\cite{ross_1960_square}. They introduced a structural characterization of the square of a graph in terms of clique intersections, multicliqual 
vertices, and so called "neighborly" properties, proving that a graph is the square of a tree if and only if
\begin{enumerate}[label=\roman*)]
  \item every vertex is "neighborly" - is not isolated, 
  \item no two cliques intersect in more than two vertices, and
  \item the number of clique pairs intersecting in two vertices is exactly one less than the total number of cliques. 
\end{enumerate}
In addition to these properties, there must exist a one-to-one correspondence between maximal cliques and multicliqual vertices where a vertex is
called multicliqual if it belongs to at least two maximal cliques. If a multicliqual vertex $b$ belongs to exactly $k$ maximal cliques, 
then the clique corresponding to $b$ contains exactly $k$ multicliqual vertices. Furthermore, if two maximal cliques intersect in
exactly one vertex $b$, then there exists a third maximal clique which contains $b$ and shares exactly one additional vertex 
with each of the original two cliques. Using these properties, they constructed an explicit algorithm that recovers the underlying tree 
by associating multicliqual vertices with maximal cliques and forming a tree from pairwise clique intersections.

Since then, the problem of finding the square root of a graph has extended beyond sociology and is a well-studied field in both mathematics and 
computer-science, with direct applications in distributed computing (e.g., see \cite{peres_2023_distributed}), network design and computational 
biology. Research in this area has primarily focused around an undirected graph $H$ and its square $G = H^2$ which is obtained by adding an edge 
between every pair of vertices whose distance in $H$ is at most two. 

This two-step undirected graph square has been studied extensively with much of the focus being on the complexity of the decision problem - 
given a specific square graph $G$ does there exist a root graph $H$ such that $H^2=G$. \citeauthor{motwani_sudan_1994_roots} showed that graph 
squares are NP-complete in general and hence subsequent studies in the area have restricted the decision problem in two ways, either by considering
various properties on the root graph or by considering properties on the square graph. 

\section{Aims}

In contrast, rather than studying the construction problem directly, in this work we analyze the relationship between the adjacency matrix $A_H$ 
of the root and the adjacency matrix $A_G$ of the square, and show that the problem can be completely categorized as an integer program. We then 
examine the relaxed linear program which can be evaluated in polynomial time using interior point method, and integrate existing literature on special root 
structures as a stepping-stone towards generating a true polynomial time algorithm.  

We begin by introducing the comprehensive literature surrounding the subject and provide a general overview of the research direction in this area 
and the major advances that have been made over the years. We then standardize the definitions that will be used for the rest of this thesis and 
introduce fundamental theory on the three areas of complexity theory, graph theory and mathematical optimization that are used in the various 
algorithmic approaches to the decision problem. 

The first algorithm we explore is the polynomial time algorithm given by \citeauthor{ross_1960_square} for constructing a tree root from a given 
square. We then explore the proof used by \citeauthor{motwani_sudan_1994_roots} in showing that 


\section{Literature}
Focusing first on special undirected graph classes, the aforementioned work by \citeauthor{ross_1960_square} in 1960 gives an algorithm for 
constructing a tree root by using characterizations of the square graph. As part of this work they proved tree square roots are unique up to 
isomorphism. This was the first major contribution to the study of graph roots and provided fundamental structural characteristics of 
square graphs, laying the groundwork for later research on graph root problems and their applications \cite{ross_1960_square}. 

\citeauthor{fleischner_1974_square} in 1974 showed that as long as a graph is 2-connected its square graph will have a Hamiltonian cycle. This 
established a broad condition guaranteeing Hamiltonicity in graph squares which is used in many subsequent works \cite{fleischner_1974_square}. 

Lau and Corneil in 2004 provided an algorithm for constructing square roots from an interval square graph in polynomial time, and showed that the 
decision problem for finding roots of chordal graphs, roots of split graphs, and chordal roots of graphs are NP-complete
\cite{lau_corneil_2004_chordal_root}. Furthermore, Lau considered the construction of square roots from a bipartite graph in linear 
time, and showed that the existence of cube roots in bipartite square is NP-complete \cite{lau_2006_bipartite_roots}. 

Le and Tuy in 2007 generalized tree-root results further by giving a linear time recognition of roots of block graphs \cite{le_tuy_2007_block_square}. 

Extending to the broader generalizing of graph squares, Mukhopadhyay in 1967 provided necessary and sufficient conditions for a connected graph 
to have a generalized square root \cite{mukhopadhyay_1967_root}, whilst Geller in 1970 likewise found necessary and sufficient conditions for a 
digraph to have a square root \cite{geller_1970_root}. 

In 1994 Motwani and Sudan established the complexity of the square problem for the generalized undirected case and
showed that the problem of determining if a given graph has a square root is a NP-complete problem via a polynomial reduction to 3-CNF
Not-All-Equal SAT \cite{motwani_sudan_1994_roots}.

Their work led to joint work by Farzad, Le, Tuy, Lau and Karimi in 2009 and 2012, in which they established a partial dichotomy theorem for 
the complexity of determining if a given undirected graph of girth $n$ has a square root \cite{farzad_2009_roots}
which was informally completed under \cite{farzad_karimi_2012_root}. 
In 2013, Cochefert and others gave an algorithm for constructing a square root with minimal number of edges from undirected graphs with max degree 
$\leq 6$ in polynomial time and an algorithm for constructing maximal edge square roots in exponential time \cite{cochefert_et_al_2013_sparse_square_roots}. 

In recent years, Duffy provided a full classification for the undirected square of oriented graphs
satisfying max degree $\leq 3$ and girth $\geq 4$ \cite{duffy_2025_oriented_root}. 


\section{Definitions}

\subsection{Complexity Theory}

Complexity theory studies the classification of computational problems and explores the amount of resources required to solve 
certain problems. The critical resource that is measured is running time and is standardized across various different problems by using

\begin{definition}[Big-$O$ Notation] \label{def::big_O_notation}
  Let $f,g:\mathbb{N} \rightarrow \mathbb{R}_{\geq 0}$. The \textbf{Big-$O$ Notation} $f(n)=O(g(n))$ is written 
  if there exist constants $c>0$ and $n_0\in\mathbb{N}$ such that for all $n\geq n_0$
  \[
  f(n)\leq c g(n). 
  \]
\end{definition}

Using this notation, can classify algorithms based on their runtime. 

\begin{definition}[Polynomial-Time Algorithm] \label{def::polynomial_time_algorithm}
  An algorithm is said to run in \textbf{polynomial time} if there exists a constant $c \ge 0$ such that its
  running time on inputs of size $n$ is $O(n^c)$.
\end{definition}

Polynomial-time algorithms are important because they are generally regarded as efficient and computationally ``simple". Consequently, 
the class of problems that admit polynomial-time algorithms forms the basis for the complexity class $\mathbf{P}$.

\begin{definition}[Decision Problem] \label{def::decision_problem}
  A \textbf{decision problem} is a computational problem whose output is either \textbf{Yes} or \textbf{No} for 
  every valid input instance. 
\end{definition}

\begin{definition}[Class $\mathbf{P}$] \label{def::class_p}
  The complexity class $\mathbf{P}$ consists of all decision problems that can be solved by a deterministic 
  algorithm in polynomial time.
\end{definition}

Unfortunately many important decision problems are not known to admit such polynomial-time algorithms and hence need additional classification
on their ``hardness''. For these problems, despite not have an efficient algorithm for finding a solution, a proposed solution can be
verified efficiently.

\begin{definition}[Certificate] \label{def::certificate}
  Let $L$ be a decision problem. A \textbf{certificate} for an instance $x$ of $L$ is a finite string $c$ that 
  provides evidence that $x$ is a \textbf{Yes}-instance of $L$. A certificate is said to be 
  \textbf{polynomially verifiable} if there exists a deterministic algorithm that for given input $(x,c)$, can 
  determines whether $c$ is a valid certificate for $x$ in time polynomial in $|x|$.
\end{definition}

\begin{definition}[Class $\mathbf{NP}$] \label{def::class_np}
  The complexity class $\mathbf{NP}$ consists of all decision problems for which every \textbf{Yes}-instance admits 
  a certificate that is polynomially verifiable. 
\end{definition}

Since every decision problem in $\mathbf{P}$ can be solved in polynomial time, a solution can also be verified in polynomial time,
\[
\mathbf{P} \subseteq \mathbf{NP}.
\]
Whether this inclusion is strict remains one of the most important open questions in theoretical computer science - the
$\mathbf{P}$ vs $\mathbf{NP}$ problem. As $\mathbf{NP}$ is a large space, the relative difficulty of problems in $\mathbf{NP}$
is compared by using reduction. 

\begin{definition}[Polynomial-Time Reduction] \label{def::polynomial_time_reduction}
  Let $A$ and $B$ be decision problems. A \textbf{polynomial-time reduction} from $A$ to $B$, denoted $A \leq_p B$,
  is a polynomial-time computable function $f$ such that, for every instance $x$ of $A$,
  \[
  x \in A \Longleftrightarrow f(x) \in B.
  \]
\end{definition}


\begin{definition}[Class $\mathbf{NP}$-Hard] \label{def::class_np_hard}
  A decision problem $B$ is \textbf{$\mathbf{NP}$-hard} if every problem $A \in \mathbf{NP}$ satisfies $A \leq_p B$. 
  The collection of all $\mathbf{NP}$-hard problems is denoted by the class \textbf{$\mathbf{NP}$-hard}.
\end{definition}

\begin{definition}[Class $\mathbf{NP}$-Complete] \label{def::class_np_complete}
  A decision problem $B$ is \textbf{$\mathbf{NP}$-complete} if $B \in \mathbf{NP}$, and $B$ is $\mathbf{NP}$-hard.
  The collection of all $\mathbf{NP}$-complete problems is denoted by the class \textbf{$\mathbf{NP}$-complete}.
\end{definition}


\subsection{Graph Theory}

\begin{definition}[Graph] \label{def::graph}
  
\end{definition}

\begin{definition}[Undirected Graph] \label{def::undirected_graph}
  
\end{definition}

\begin{definition}[Directed Graph] \label{def::directed_graph}
  
\end{definition}

\begin{definition}[Graph Distance] \label{def::graph_distance}
  
\end{definition}

\begin{definition}[Graph Root] \label{def::graph_root}
  
\end{definition}

\begin{definition}[Graph Square] \label{def::grapH_square}
  
\end{definition}

\begin{definition}[Adjacency Matrix] \label{def::adjacency_matrix}
  
\end{definition}

\subsection{Mathematical Optimization}

\begin{definition}[Integer Program] \label{def::integer_program}
  
\end{definition}

\begin{definition}[Linear Program] \label{def::linear_program}
  
\end{definition}


Let the \emph{graph root} $H = (V, E_H)$ be a finite, simple, undirected graph, where $V$ is the vertex set with $|V| = n$, 
and $E_H \subseteq \{\{u,v\} \mid u,v \in V,\; u \neq v\}$ is the edge set. Denote by $d_H(u,v)$ the 
\emph{graph distance} defined as the number of edges of the shortest path 
between vertices $u$ and $v$ in $H$. If no such path exists, $d_H(u,v) = \infty$.

For vertices $u,v\in V$, the \emph{graph distance}
$d_H(u,v)$ is the length of a shortest path between
$u$ and $v$ in $H$. If no such path exists, then
$d_H(u,v)=\infty$.

The \emph{square} of the graph $H$, denoted by $H^2 = (V, E_{H^2})$ is the graph on the same vertex set with edge set
\begin{equation} \label{eq::fund_def_square}
E_{H^2} = \{\, uv \mid u \neq v,\; 1 \leq d_H(u,v) \leq 2 \,\}.
\end{equation}
That is, two distinct vertices are adjacent in $H^2$ if and only if they are connected by a path of length one or two in $H$. 
This is the standard definition of the square of a graph.

Using an alternative construction, let the \emph{graph square} $G = (V, E_G)$ be an undirected graph defined on 
the same vertex set $V$. Let $A_H \in \{0,1\}^{n \times n}$ denote the adjacency matrix of $H$, where
\[
(A_H)_{ij} =
\begin{cases}
1 ,& \text{if } ij \in E_H, \\
0 ,& \text{otherwise},
\end{cases}
\]
and let $D_H \in \mathbb{R}^{n \times n}$ denote the degree matrix of $H$, where
\[
(D_H)_{ij} =
\begin{cases}
\deg_H(i) ,& \text{if } i = j, \\
0 ,& \text{otherwise}.
\end{cases}
\]
Let $\mathbbm{1}_{\geq 1} : \mathbb{R}^{n \times n} \to \{0,1\}^{n \times n}$ denote the 
entrywise \emph{indicator function}, defined by
\[
\mathbbm{1}_{\geq 1}(M)_{ij} =
\begin{cases}
1 ,& \text{if } M_{ij} \geq 1, \\
0 ,& \text{otherwise}.
\end{cases}
\]
Define the adjacency matrix of $G$ using the adjacency and degree matrices of $H$
\begin{equation} \label{eq::fund_adj_eq_undirected}
A_G = \mathbbm{1}_{\geq 1}\!\big((A_H)^2 - D_H + A_H\big).
\end{equation}

\chapter{Algorithmic Approaches} \label{chap::algorithmic_approaches}

In this chapter we explore different methodologies that have been used to show certain properties of the graph root and 
the graph square. For many restricted cases which are NP, a polynomial time reduction is given and incases which are
polynomial, a polynomial time algorithm is given.  


We then examine the relaxed linear program which can be evaluated in polynomial time using interior point method, and 
integrate existing literature on special root structures as a stepping-stone towards generating a true polynomial time 
algorithm. 

\section{Polynomial Time Algorithm} 

\citeauthor{ross_1960_square} give an algorithm for constructing a tree $T$ when given its square graph $G = T^2$ 
using the following characterization.  
\begin{itemize}
    \item If $G = K_p$ is complete, then $T$ must be a tree of depth 1 (result of \hyperref[lemma:1]{Lemma 1})
    \item If $G \neq K_p$ 
    \begin{enumerate}[label=(\roman*)]
        \item\label{characterization:1} Every vertex of $G$ forms a clique with its neighbours and $G$ is connected (result of \hyperref[lemma:2]{Lemma 2})
        \item\label{characterization:2} If two cliques intersect at only one point $v$ then there is a third clique with which they share $v$ and exactly
                one other point (result of \hyperref[lemma:3]{Lemma 3} and \hyperref[lemma:5]{Lemma 5})
        \item\label{characterization:3} A clique $C(v)$ in $G$ corresponding to $v$ in $T$ contains exactly as many points in multiple cliques
            as the number of cliques which include $v$ (result of \hyperref[lemma:3]{Lemma 3} and \hyperref[lemma:4]{Lemma 4})
        \item\label{characterization:4} Every pair of cliques intersect at at most two points (result of \hyperref[lemma:5]{Lemma 5})
        \item\label{characterization:5} The number of pairs of cliques that meet at two points is one less than the number of cliques (result of 
            cliques formed from non-leaf nodes forming a tree structure)
    \end{enumerate}
\end{itemize}

\subsubsection*{Lemma 1}\label{lemma:1}
$G$ is complete iff $T$ has depth 1. 
\begin{proof}
$\left(\leftarrow\right)$ Let $T$ be a tree with depth 1, hence all vertices in $T$ are at most 2-distance away from 
any other vertex since the root is a non-leaf vertex adjacent to all vertices in $T$ $\implies$ $G = T^2$ is complete since 
edges exist between all vertices

$\left(\rightarrow\right)$ Let $G = T^2$ be complete, then all vertices must have distance at most two in $T$
and $T$ has depth at most 2. Since $T$ is a tree, there are no cycles and hence there cannot be depth 2 since this 
would lead to a path of length 3 from the vertex in depth 2 to an alternate vertex in depth 1. Hence $T$ has exactly 
depth 1. 
\end{proof}

\subsubsection*{Lemma 2}\label{lemma:2}
Every vertex of $G$ forms a clique with its neighbours and $G$ is connected
\begin{proof}
For a vertex $v \in T$, its neighbours in $T$ are all adjacent edge in $G$ by construction 
hence $\left\{ v \right\} \cup \Gamma \left(v\right)$ is a clique in $G$.
$T \subseteq G$ and since $T$ is a tree which is connected, $G$ must also be connected. 
\end{proof}

\subsubsection*{Lemma 3}\label{lemma:3}
Every maximal clique of $G$ corresponds to a closed neighbourhood of a non-leaf vertex and 
every closed neighbourhood of a non-leaf vertex in $T$ form a clique in $G$

\begin{proof}
A set of leaves in $T$ cannot form a clique in $G$ without having at least one non-leaf since the distance between 
every leaf is at least greater than 2 unless they share a common neighbour. 

Consider a non-leaf vertex $v \in T$, since all vertices in neighbourhood $\Gamma\left(v\right)$ have 
distance at most 2 in $G$, they will form a clique in $G$. This clique is maximal since all neighbours of $v$ are already 
part of the clique and no additional vertex can be added. 
\end{proof}

\subsubsection*{Lemma 4}\label{lemma:4}
If $G$ has more than one maximal clique, then vertex $v$ is in exactly one maximal clique in $G$ iff $v$ is a leaf in $T$

\begin{proof}
$\left(\leftarrow\right)$ Suppose vertices $v, w$ to be in two cliques, there must be a neighbour $x \in \Gamma\left(v\right)$ 
that is not adjacent to $w$. If $x$ is not a leaf in $T$, will form a third clique in $G$ and being 2-distance away from
$w$ will also be in the clique formed from $w$. 
Hence only if this adjacent vertex is a leaf will it be in only one clique $G$.  

$\left(\rightarrow\right)$ Suppose $v$ is a leaf in $T$ whereby $v$ has exactly one neighbour $w$ with is a non-leaf. 
The only maximal clique in $G$ that includes $v$ must also include $w$ and the neighbourhood $\Gamma\left(w\right)$.
Since $v$ has no other neighbours it cannot appear in any other clique and hence is in exactly one clique.  
\end{proof}

\subsubsection*{Lemma 5}\label{lemma:5}
Two non-leaves of $T$ are adjacent in $T$ iff their neighbourhoods are cliques of $G$ that intersect in exactly two points

\begin{proof}
$\left(\leftarrow\right)$ Supposes cliques in $G$ correspond to the neighbourhoods of two non-leaf neighbours $v, w$ 
which intersect at exactly $v, w$. Then $v, w$ must be adjacent, otherwise if they were instead 2-distance away then
they would share a common neighbour $u$ which is also in both cliques. 

$\left(\rightarrow\right)$ Let $v, w$ be adjacent vertices that are non-leafs in $T$. Both will form cliques in $G$ and be in each clique since they 
are neighbours. Assume to the contrary that $\exists$ vertex $k$ also in both cliques. $k$ must be adjacent to both 
$v, w$ in $T$ but this creates a cycle since there is now two different paths between $v, w$. Hence $v, w$ are the 
only two vertices that intersect in the two clique neighbourhoods.   
\end{proof}

\subsubsection*{Lemma 6}\label{lemma:6}
$G$ is a connected graph where the removal of any edge does not result in a disconnected graph

\begin{proof}
By construction $\exists$ two paths between every pair of vertices $v, w \in G$ where one of the vertices $v$ must
be non-leaf. Remove the existing $vw$ edge in $T \subset G$. There still exists a path by taking the 2-distance edge 
to the neighbour $u$ of $v$ and then going from $u to w$ ($vu, uw$), hence the graph $G$ is still connected.
\end{proof}

Combining this gives the following 

\subsection*{Algorithm}
\begin{algorithm}
\begin{algorithmic}[1]
\State \textbf{Input:} A graph $G$ 
\State \textbf{Output:} A tree $T$ where $G = T^2$
\vspace{0.5em}

\Statex \textbf{Step 1}
\If{$G = K_p$ is complete}
    \State \textbf{Return} tree $T$ with depth $1$ where root can be any vertex
\Else
    \State \textbf{Continue} 
\EndIf
\vspace{0.5em}

\Statex \textbf{Step 2}
\State Identify all maximal cliques $C_1, \dots, C_n$ of $G$ using an algorithm such as Bron–Kerbosch where $n > 1$ 
using \hyperref[characterization:1]{characterization (i)}
\vspace{0.5em}

\Statex \textbf{Step 3}
\State Let $b_1, \dots, b_n$ be multicliqual points corresponding to $C_1, \dots, C_n$ formed from 
\hyperref[characterization:3]{characterization (iii)} 
\State Construct a graph $S$ with vertex set $\{b_1, \dots, b_n\}$
\For{each pair $(C_i, C_j)$ satisfying $|C_i \cap C_j| = 2$  }
    \State Add edge between $b_i$ and $b_j$ in $S$
\EndFor
\State By \hyperref[characterization:5]{characterization (v)}, $S$ is a tree

\vspace{0.5em}
\Statex \textbf{Step 4}
\For{each clique $C_i$}
    \State Let $n_i$ be the number of unicliqual points in $C_i$
    \State Attach $n_i$ leaf nodes to $b_i$ in $S$
\EndFor
\State The resulting tree $T$ has the same number of nodes as $G$ and satisfies $T^2 = G$
\end{algorithmic}
\end{algorithm}


\section{NP-Completeness Proofs}

\section{Integer Program Foundation}

\subsection{Background}

\begin{lemma} \label{lemma::G_equal_to_H^2}
The graph $G$ defined by \eqref{eq::fund_adj_eq_undirected} is equivalent to the graph square $H^2$
 defined in \eqref{eq::fund_def_square} 
\begin{equation} \label{eq::fund_def_square_expanded}
 ij \in E_{H^2} \iff 1 \leq d_H(i,j) \leq 2. 
\end{equation}
\end{lemma}
\begin{proof} 
\noindent
For $i \neq j$, since $D_H$ is diagonal, have $(D_H)_{ij} = 0$ and hence
\[
\big((A_H)^2 - D_H + A_H\big)_{ij} = (A_H)^2_{ij} + (A_H)_{ij}.
\]

($\Rightarrow$) Suppose $ij \in E_G$ such that $(A_H)^2_{ij} + (A_H)_{ij} \geq 1$.
This implies that at least on of $(A_H)_{ij} = 1$ or $(A_H)^2_{ij} \geq 1$ must be true. If $(A_H)_{ij} = 1$, then $d_H(i,j) = 1$.  
If $(A_H)^2_{ij} \geq 1$, then there exists a vertex $k$ such that $ik, kj \in E_H$, and hence $d_H(i,j) \leq 2$. Therefore 
$1 \leq d_H(i,j) \leq 2$, and thus $ij \in E_{H^2}$.

($\Leftarrow$) Suppose $ij \in E_{H^2}$ such that $1 \leq d_H(i,j) \leq 2$. 
$D_H$ is diagonal such that $(D_H)_{ij} = 0$. If $d_H(i,j) = 1$, then $(A_H)_{ij} = 1$.  
If $d_H(i,j) = 2$, then there exists a vertex $k$ such that $ik, kj \in E_H$, and hence $(A_H)^2_{ij} \geq 1$.
In either case, $(A_H)^2_{ij} + (A_H)_{ij} \geq 1$, so $(A_G)_{ij} = 1$, and hence $ij \in E_G$.

For $i = j$, since all graphs are simple, have $(E_{\cdot})_{ii} = 0$, $(A_{\cdot})_{ii} = 0$ and hence
\[
\big((A_H)^2 - D_H + A_H\big)_{ii} = (A_H)^2_{ii} - (D_H)_{ii}.
\]
$(A_H)^2_{ii} - (D_H)_{ii} = 0$ is always true since the degree matrix of $H$ always has diagonal 
values $\displaystyle (D_H)_{ii} = \sum_k x_{ik}$ and each edge from vertex $i$ to its neighbour in $H$ 
will always create a loop to itself in the square such that $\displaystyle (A_H)^2_{ii} = 
\sum_{k = 1} x_{ik}x_{ki} = \sum_{k = 1} x_{ik}^2$ by the symmetry of adjacency matrix $A_H$. Since $x_{ik} \in \{0, 1\}$,
the square of the boolean is identical to the boolean $x_ik$ such that $\displaystyle (A_H)^2_{ii} = \sum_k x_{ik} = (D_H)_{ii}$
and $(A_G)_{ii} = \mathbbm{1}_{\geq 1}\!\big((A_H)^2 - D_H + A_H\big)$ as required. 

Combining both gives that $E_G = E_{H^2}$, and so $G = H^2$.  

\end{proof}
\noindent
The formulation of the adjacency matrix of the graph square in Equation~\eqref{eq::fund_adj_eq_undirected} is very powerful. 
However, the presence of the indicator function introduces a nonlinearity that prevents the problem from being expressed 
as a purely quadratic form, thereby limiting the direct applicability of existing results from the quadratic optimization literature.

\begin{lemma} \label{lemma::no_indicator_girth_5}
  If there are no $K_3$ triangle subgraphs or any $C_4$ square subgraphs in the graph root such that $H$ has girth of length $5$,
  then  Equation~\eqref{eq::fund_adj_eq_undirected} can be simplified to become 
  \begin{equation} \label{eq::fund_adj_eq_undirected_no_indicator}
    A_G = \big((A_H)^2 - D_H + A_H\big).
  \end{equation}
\end{lemma}

\begin{proof}
  Since the graph root has girth of length $5$ there are no cycles of length $3$ or $4$ in $H$. 
  Thus for every pair of vertices $i \neq j \in V$ there is either no cycle that connects $i,j$ or there
  is a cycle of at least length $5$ between them. In all cases $(D_H)_{ij} = 0$. 
  If there is no cycle, then there is at most one path between $i,j$
  \begin{enumerate}[label=\roman*)]
    \item If the path is of length $1$, then $(A_H)_{ij} = 1, (A_H)^2_{ij} = 0 \implies (A_G)_{ij} = 0 - 0 + 1 = 1$.
    \item If the path is of length $2$, then $(A_H)_{ij} = 0, (A_H)^2_{ij} = 1 \implies (A_G)_{ij} = 1 - 0 - 0 = 1$.
    \item If the path is of length $\geq 3$ or a path does not exist, then $(A_H)_{ij} = 0, (A_H)^2_{ij} = 0 \implies (A_G)_{ij} = 0$.
  \end{enumerate}
  If there is a cycle of length $\geq 5$ between $i,j$ such that there are two paths between $i,j$
  \begin{enumerate}[label=\roman*)]
    \item If one path is of length $1$, then the other is of length $\geq 4$ and hence $(A_H)_{ij} = 1, (A_H)^2_{ij} = 0 \implies (A_G)_{ij} = 0 - 0 + 1 = 1$.
    \item If one path is of length $2$, then the other is of length $\geq 3$ and hence $(A_H)_{ij} = 0, (A_H)^2_{ij} = 1 \implies (A_G)_{ij} = 1 - 0 - 0 = 1$.
    \item If the paths are both of length $\geq 3$,  then $(A_H)_{ij} = 0, (A_H)^2_{ij} = 0 \implies (A_G)_{ij} = 0$. 
  \end{enumerate}
  If $i = j$, then \((A_H)_{ii} = 0\), \((D_H)_{ii} = \deg_H(i)\), and \((A_H)^2_{ii} = \deg_H(i)\). Thus,
  \[
  (A_H)^2_{ii} - (D_H)_{ii} + (A_H)_{ii}
  = \deg_H(i) - \deg_H(i) + 0 = 0 = (A_G)_{ii}.
  \]
\end{proof}

\section{Linear Program Relaxation}


\newpage





\printbibliography




\end{document}
